\textbf{Phần 1. Trắc nghiệm lựa chọn}
\setcounterpageref{ex}{0}

\begin{ex} %[5D4N23]
Khi nhân một số thập phân với $100$, ta chỉ việc chuyển dấu phẩy của số đó sang bên phải
\choice
{$1$ chữ số}
{\True $2$ chữ số}
{$3$ chữ số}
{$4$ chữ số}
\loigiai{
Muốn nhân một số thập phân với $100$, ta chuyển dấu phẩy của số đó sang bên phải hai chữ số.\\
\textbf{Đáp án của câu hỏi là B}
}
\end{ex}

\begin{ex} %[5D4N23]
Khi chia một số thập phân cho $0{,}1$, ta chỉ việc
\choice
{Chuyển dấu phẩy sang bên trái $1$ chữ số}
{Chuyển dấu phẩy sang bên trái $2$ chữ số}
{\True Chuyển dấu phẩy sang bên phải $1$ chữ số}
{Chuyển dấu phẩy sang bên phải $2$ chữ số}
\loigiai{
Chia cho $0{,}1$ tương đương với nhân với $10$, ta chuyển dấu phẩy sang bên phải một chữ số.\\
\textbf{Đáp án của câu hỏi là C}
}
\end{ex}

\begin{ex} %[5D4H23]
Kết quả của phép tính $2{,}08 \times 1\,000$ là
\choice
{$20{,}8$}
{\True $2\,080$}
{$208$}
{$20\,800$}
\loigiai{
Chuyển dấu phẩy sang phải ba chữ số: $2{,}08 \times 1\,000 = 2\,080$.\\
\textbf{Đáp án của câu hỏi là B}
}
\end{ex}

\begin{ex} %[5D4H23]
Kết quả của phép tính $15{,}31 \times 0{,}001$ là
\choice
{$15\,310$}
{$1{,}531$}
{$0{,}1531$}
{\True $0{,}01531$}
\loigiai{
Chuyển dấu phẩy sang trái ba chữ số: $15{,}31 \times 0{,}001 = 0{,}01531$.\\
\textbf{Đáp án của câu hỏi là D}
}
\end{ex}

\begin{ex} %[5D4H23]
Kết quả của phép tính $0{,}8 : 100$ là
\choice
{$80$}
{\True $0{,}008$}
{$0{,}08$}
{$0{,}0008$}
\loigiai{
Chuyển dấu phẩy sang trái hai chữ số: $0{,}8 : 100 = 0{,}008$.\\
\textbf{Đáp án của câu hỏi là B}
}
\end{ex}

\begin{ex} %[5D4H23]
Kết quả của phép tính $6{,}3 : 0{,}01$ là
\choice
{\True $630$}
{$0{,}063$}
{$0{,}63$}
{$63$}
\loigiai{
Chuyển dấu phẩy sang phải hai chữ số: $6{,}3 : 0{,}01 = 630$.\\
\textbf{Đáp án của câu hỏi là A}
}
\end{ex}

\begin{ex} %[5D4H23]
Đổi $23{,}45\text{ kg}$ ra gam ta được kết quả là
\choice
{$234{,}5\text{ g}$}
{$2\,345\text{ g}$}
{$234\,500\text{ g}$}
{\True $23\,450\text{ g}$}
\loigiai{
Vì $1\text{ kg} = 1\,000\text{ g}$ nên $23{,}45\text{ kg} = 23{,}45 \times 1\,000 = 23\,450\text{ g}$.\\
\textbf{Đáp án của câu hỏi là D}
}
\end{ex}

\begin{ex} %[5D4H23]
Đổi $12{,}6\text{ cm}$ ra mét ta được kết quả là
\choice
{\True $0{,}126\text{ m}$}
{$1{,}26\text{ m}$}
{$0{,}0126\text{ m}$}
{$126\text{ m}$}
\loigiai{
Vì $1\text{ m} = 100\text{ cm}$ nên $12{,}6\text{ cm} = 12{,}6 : 100 = 0{,}126\text{ m}$.\\
\textbf{Đáp án của câu hỏi là A}
}
\end{ex}

\begin{ex} %[5D4H23]
Mỗi tờ giấy dày $0{,}1\text{ mm}$. Cả chồng giấy dày $36{,}5\text{ mm}$. Số tờ giấy trong chồng đó là
\choice
{$36$ tờ}
{$3\,650$ tờ}
{\True $365$ tờ}
{$3{,}65$ tờ}
\loigiai{
Số tờ giấy là: $36{,}5 : 0{,}1 = 365$ (tờ).\\
\textbf{Đáp án của câu hỏi là C}
}
\end{ex}

\begin{ex} %[5D4V23]
Một đàn gấu Bắc Cực có $10$ chú gấu con, mỗi chú gấu con ăn hết $4{,}5\text{ kg}$ cá mỗi ngày. Hỏi cả $10$ chú gấu con ăn hết bao nhiêu ki-lô-gam cá mỗi ngày?
\choice
{\True $45\text{ kg}$}
{$450\text{ kg}$}
{$4{,}5\text{ kg}$}
{$40{,}5\text{ kg}$}
\loigiai{
Cả $10$ chú gấu con ăn hết: $4{,}5 \times 10 = 45\text{ (kg cá)}$.\\
\textbf{Đáp án của câu hỏi là A}
}
\end{ex}

\begin{ex} %[5D4V23]
Chiếc máy xúc lớn nhất thế giới xúc $10$ lần được $937{,}8\text{ tấn}$ quặng. Trung bình mỗi lần máy xúc đó xúc được số tấn quặng là
\choice
{$9\,378\text{ tấn}$}
{$9{,}378\text{ tấn}$}
{\True $93{,}78\text{ tấn}$}
{$937{,}8\text{ tấn}$}
\loigiai{
Mỗi lần máy xúc xúc được: $937{,}8 : 10 = 93{,}78\text{ (tấn)}$.\\
\textbf{Đáp án của câu hỏi là C}
}
\end{ex}

\begin{ex} %[5D4V23]
Chia đều $0{,}3\text{ l}$ giấm vào $10$ ống nghiệm như nhau. Mỗi ống nghiệm chứa số lít giấm là
\choice
{$3\text{ l}$}
{$0{,}3\text{ l}$}
{$0{,}003\text{ l}$}
{\True $0{,}03\text{ l}$}
\loigiai{
Mỗi ống nghiệm chứa: $0{,}3 : 10 = 0{,}03\text{ (l)}$.\\
\textbf{Đáp án của câu hỏi là D}
}
\end{ex}

\textbf{Phần 2. Trắc nghiệm đúng sai}
\setcounterpageref{ex}{0}

\begin{ex} %[5D4H23]
Xét tính đúng sai của các khẳng định sau:
\choiceTF
{\True $79{,}6 \times 0{,}01 = 0{,}796$.}
{$58 : 0{,}001 = 0{,}058$.}
{\True $6{,}89\text{ l} = 6\,890\text{ ml}$.}
{$3{,}7\text{ m} = 0{,}037\text{ km}$.}
\loigiai{
\begin{itemize}
    \item $79{,}6 \times 0{,}01 = 0{,}796$.
    \item $58 : 0{,}001 = 58 \times 1\,000 = 58\,000 \neq 0{,}058$.
    \item $6{,}89\text{ l} = 6{,}89 \times 1\,000 = 6\,890\text{ ml}$.
    \item $3{,}7\text{ m} = 3{,}7 : 1\,000 = 0{,}0037\text{ km} \neq 0{,}037\text{ km}$.
\end{itemize}
\begin{itemize}
    \item a) \textbf{Đúng}.
    \item b) \textbf{Sai}: Kết quả là $58\,000$.
    \item c) \textbf{Đúng}.
    \item d) \textbf{Sai}: $3{,}7\text{ m} = 0{,}0037\text{ km}$.
\end{itemize}
}
\end{ex}

\begin{ex} %[5D4V23]
Cho dãy biến đổi tính toán: $A = 10{,}8 \xrightarrow{: 100} B \xrightarrow{: 0{,}01} C \xrightarrow{\times 0{,}01} D \xrightarrow{\times 100} E$. Xét tính đúng sai của các khẳng định sau:
\choiceTF
{\True Giá trị của $B$ là $0{,}108$.}
{Giá trị của $C$ là $1{,}08$.}
{\True Giá trị của $D$ bằng giá trị của $B$.}
{\True Giá trị cuối cùng $E$ bằng $10{,}8$ (bằng với số ban đầu $A$).}
\loigiai{
\begin{itemize}
    \item $B = 10{,}8 : 100 = 0{,}108$.
    \item $C = 0{,}108 : 0{,}01 = 10{,}8 \neq 1{,}08$.
    \item $D = 10{,}8 \times 0{,}01 = 0{,}108 = B$.
    \item $E = 0{,}108 \times 100 = 10{,}8 = A$.
\end{itemize}
\begin{itemize}
    \item a) \textbf{Đúng}.
    \item b) \textbf{Sai}: $C = 10{,}8$.
    \item c) \textbf{Đúng}.
    \item d) \textbf{Đúng}.
\end{itemize}
}
\end{ex}

\textbf{Phần 3. Trả lời ngắn}
\setcounterpageref{ex}{0}

\begin{ex} %[5D4H23]
Tính kết quả phép tính: $0{,}94 \times 100$.
\shortans{94}
\loigiai{
\textbf{Đáp số:} $94$\\
$0{,}94 \times 100 = 94$.
}
\end{ex}

\begin{ex} %[5D4H23]
Tính kết quả phép tính: $8{,}2 \times 0{,}01$. (Nhập kết quả dưới dạng số thập phân).
\shortans{0,082}
\loigiai{
\textbf{Đáp số:} $0{,}082$\\
$8{,}2 \times 0{,}01 = 0{,}082$.
}
\end{ex}

\begin{ex} %[5D4H23]
Tính kết quả phép tính: $341 : 1\,000$. (Nhập kết quả dưới dạng số thập phân).
\shortans{0,341}
\loigiai{
\textbf{Đáp số:} $0{,}341$\\
$341 : 1\,000 = 0{,}341$.
}
\end{ex}

\begin{ex} %[5D4H23]
Đổi $23{,}4\text{ tạ}$ ra yến.
\shortans{234}
\loigiai{
\textbf{Đáp số:} $234$\\
$23{,}4\text{ tạ} = 23{,}4 \times 10 = 234\text{ yến}$.
}
\end{ex}

\begin{ex} %[5D4V23]
Một tấm bạt hình chữ nhật có chiều dài $27{,}86\text{ m}$ và chiều rộng $10\text{ m}$. Tính diện tích tấm bạt đó theo đơn vị mét vuông. (Nhập kết quả dưới dạng số thập phân).
\shortans{278,6}
\loigiai{
\textbf{Đáp số:} $278{,}6$\\
$27{,}86 \times 10 = 278{,}6\text{ (m}^2\text{)}$.
}
\end{ex}

\begin{ex} %[5D4V23]
Một thanh sắt dài $1\text{ m}$ nặng $15{,}23\text{ kg}$. Hỏi thanh sắt cùng loại dài $0{,}1\text{ m}$ nặng bao nhiêu ki-lô-gam? (Nhập kết quả dưới dạng số thập phân).
\shortans{1,523}
\loigiai{
\textbf{Đáp số:} $1{,}523$\\
$15{,}23 \times 0{,}1 = 1{,}523\text{ (kg)}$.
}
\end{ex}
